\section{Sonuç}
\label{AltinciBolum}

Projenin ilk dikey dilimi ve model karşılaştırması tamamlanmıştır. Veri setinin
resmî bölmeleri doğrulanmış, açılamayan dosyalar ayıklanmış ve tüm videolar için
iskelet koordinatları çıkarılmıştır. 226 sınıfta 60 epoch eğitilen iki katmanlı
BiGRU modeli yüzde 88,24 test doğruluğuna, yüzde 30 daha az parametreye sahip
1B-CNN + Transformer modeli ise yüzde 88,83 test doğruluğuna (doğrulama yüzde
89,97) ulaşarak yüzde 85 hedefini aşmıştır. Temel modelin int8 LiteRT sürümü
3,19 MB'tır ve 10 MB'lık boyut hedefinin oldukça altındadır. Python ve telefon
çıktıları arasındaki yüzde 99,5 uyum, dönüşümün model davranışını bozmadığını
göstermektedir.

Bununla birlikte iki konu henüz kapanmamıştır. Birincisi hız: model çıkarımı 48 ms
olsa da uçtan uca hat yaklaşık 9 FPS'te kalmakta ve 15 FPS hedefini
tutmamaktadır. İkincisi canlı kullanım: kaydedilmiş test verisinde yüzde 88--89
doğruluk alınırken canlı kamera denemelerinde tahmin güvenleri yüzde 3--20
aralığına düşmektedir. Sorunun MediaPipe'ın canlı görüntüdeki nokta kalitesinden
mi, ön işleme farkından mı yoksa segmentleme kurallarından mı kaynaklandığı henüz
ayrıştırılamamıştır. Aynı videoların hem Python hem telefon üzerinde çıkarılıp
nokta bazında karşılaştırılması planlanmaktadır.

Sonraki adımlarda her iki modelin üç farklı tohumla koşularak standart sapma
aralığının raporlanması ve 1B-CNN + Transformer modelinin de LiteRT int8
formatına aktarılarak telefon üzerindeki çıkarım süresinin BiGRU ile
karşılaştırılması hedeflenmektedir. Saha değerlendirmesi için ön kamera
aynalaması ve ışık değişimlerini kapsayan, en az 30 işaret ve farklı kullanıcıları
içeren bir gönüllü test protokolü hazırlanacaktır. Flutter ile Kotlin arasındaki
olay kanalı akışı, uçak modu ve üç ekranlı uygulama testleri de takip eden
kilometre taşlarında tamamlanacaktır.
